\documentclass[12pt]{article}

\usepackage[T1]{fontenc}
\usepackage[utf8]{inputenc}
\usepackage[brazil]{babel}
\usepackage{lmodern}
\usepackage{microtype}
\usepackage[a4paper,margin=2.5cm]{geometry}
\usepackage{setspace}
\usepackage{indentfirst}

\setstretch{1.15}

\title{Marginalismo e o Methodenstreit \thanks{Texto Traduzido livremente por Luiz Mario Andrade com o auxílio da ferramenta Chat GPT 5.4}}
\author{Autor do Texto}
\date{Data de Publicação}

\begin{document}

\maketitle

\section*{6 \quad Marginalismo e o \textit{Methodenstreit}}

‘[A] teoria pura da economia é uma ciência que se assemelha às ciências físico-matemáticas em todos os aspectos’.

\hspace{8cm} Walras (1954 [1874], pp. 71–2)

\section{Introdução}

A Escola Histórica e os marginalistas tiveram um ponto de partida comum: reformar a economia política para retificar o que consideravam ser as inadequações da escola clássica em geral e do ricardianismo em particular. No entanto, embora os historicistas pensassem que uma causa básica dos problemas da economia política clássica era o uso do método dedutivo abstrato ricardiano, os marginalistas consideravam esta a ferramenta mais apropriada juntamente com o utilitarismo de Bentham, que poderia ser usado como base para a reconstrução total da economia política. A Seção 2 examina como os ‘primeiros’ marginalistas, especialmente a troika composta por William Stanley Jevons (1835–82), Léon Walras (1834–1910) e Carl Menger (1840–1921), fundamentaram a abordagem em rupturas explícitas com a economia política clássica, cada um à sua maneira.

Pois também havia diferenças entre os primeiros marginalistas, tanto em substância quanto em contexto. Em particular, para Carl Menger, quando os marginalistas escreveram suas principais obras no início da década de 1870, a Escola Histórica Alemã era dominante nas universidades alemãs. Os \textit{Princípios de Economia} de Menger foram publicados em 1871 e receberam uma recepção mista, Bostaph (1978, p. 5), Caldwell (2004, pp. 35–8) e Streissler (1990, pp. 38–40). Uma crítica um tanto fria, se não hostil, foi muito provavelmente escrita por Gustav Schmoller (2004) [1873]. E isso, por sua vez, também pode ajudar a explicar por que Menger achou necessário dedicar suas energias para produzir um livro, suas \textit{Investigações}, publicado em 1883, com o objetivo singular de expor as inadequações da Escola Histórica. Este livro deu origem a um dos, se não o mais famoso dos debates metodológicos na história do pensamento econômico, o chamado \textit{Methodenstreit} ou a ‘Batalha dos Métodos’ entre Carl Menger e Gustav Schmoller. Este debate eclodiu em 1883, grassou por algum tempo, antes de desaparecer gradualmente após servir como uma pedra de toque simbólica para o compromisso relativo com a indução ou dedução, teoria ou narrativa, universalidade ou especificidade, etc.

Este debate é o foco da Seção 3. A substância do debate em torno dessas questões é de menos interesse do que seu contexto e consequências. Pois o debate teve o efeito de simbolizar e consolidar diferenças de opinião sobre as virtudes da teoria dedutiva abstrata \textit{vis-à-vis} a abordagem histórica. Ao fazer isso, ajudou a impulsionar e delinear as fronteiras entre a Escola Histórica Alemã e o marginalismo, e inaugurar a Escola Austríaca de economia, Capítulo 13.

A Seção 4 trata das questões do impacto da revolução marginalista nos desenvolvimentos que ocorreram nas ciências sociais e o destino do \textit{Methodenstreit}. Com a ascensão lenta, mas constante, da hegemonia dos seguidores da revolução marginalista, um lar não era mais encontrado dentro da própria disciplina da economia para o social e o histórico, à medida que ela se tornava cada vez mais despojada da economia política. Em vez disso, a economia política foi fragmentada, contribuindo assim para o processo de fragmentação das ciências sociais e para a formação de novas disciplinas como as conhecemos hoje. No entanto, nem toda a economia política foi retida à medida que as disciplinas de sociologia, história econômica, etc., assumiram métodos, dinâmicas e assuntos próprios. Especificamente, o conteúdo histórico da economia política de várias maneiras abriu caminho para a longeva história econômica, Capítulo 8, mas também para a efêmera economia social, Capítulo 11, se agora reemergente como a nova sociologia econômica.

Assim, quaisquer que tenham sido as causas da revolução marginalista, e quer a Escola Histórica tenha sido derrotada ou superada no \textit{Methodenstreit}, as consequências foram muito mais amplas e mistas do que simplesmente estabelecer a economia \textit{mainstream} como ela é hoje, conforme discutido nas observações finais. Exatamente o que foi descartado pelo marginalismo, e quando e como chegou ou não a ser usado em outros lugares, provou ser um processo sutil e complexo, embora o ponto final, no que diz respeito apenas à economia, tenha sido particularmente perceptível. Antes do \textit{Methodenstreit}, Jevons, por exemplo, estava muito à frente de seu tempo ao tratar a economia como um ramo das ciências naturais. O mesmo se aplica a Walras por outras razões; seu equilíbrio geral só se provou palatável uma vez que os princípios do equilíbrio parcial foram totalmente estabelecidos e aceitos, e só então estendidos à economia como um todo mais uma vez. E a promoção do marginalismo por Menger provou ser conveniente na oposição ao historicismo, mas foi menos conveniente em sua compreensão mais completa do subjetivismo (a importância tanto do significado quanto da inventividade para os indivíduos). Assim, após o \textit{Methodenstreit} e a ascensão do marginalismo, o interesse pelo método declinou, concentrando-se mais na extensão em que ele se conformava ou não às necessidades desta última em termos teóricos e técnicos, descartando tudo o que não se conformava e frequentemente avançando aos tropeços, sem levar em conta as implicações metodológicas. Inicialmente, porém, isso só poderia ser feito aceitando os argumentos dos oponentes, ainda que apenas como ressalvas em vez de concessões. Nesse aspecto, a liderança foi tomada por Alfred Marshall, especialmente no que diz respeito aos princípios, com John Neville Keynes lidando com o lado metodológico das coisas, conforme abordado no próximo capítulo.

\section{O marginalismo e o segundo cisma no pensamento econômico}

Se Marx pensava em si mesmo como alguém que fornecia uma reconstrução crítica do sistema ricardiano, Capítulo 3, o mesmo não se aplica aos marginalistas. Apesar das grandes diferenças em seus respectivos projetos, a serem discutidas mais detalhadamente abaixo e no Capítulo 7, Seção 2, o objetivo consciente dos três principais originadores da chamada revolução marginalista, Jevons, Walras e Menger, era nada menos que a reconstrução total da economia política. Embora a recepção de suas ideias tenha sido apenas gradual, seus objetivos foram plenamente realizados no longo prazo. A revolução marginalista e a subsequente emergência da economia neoclássica trouxeram uma grande ruptura no pensamento econômico. ‘As realizações combinadas de Jevons, Menger e Walras’, diz Coats (1973a, p. 38), ‘constituíram de fato um avanço intelectual significativo no desenvolvimento da análise econômica e podem ser consideradas revolucionárias em suas implicações, se não em sua novidade ou na velocidade de difusão’. Está agora bem estabelecido que a revolução marginalista não foi um evento dramático que ocorreu subitamente no início da década de 1870, mas sim uma espécie de ‘reconstrução fundamental’, Black (1973, p. 99), que foi o resultado de um processo longo e intermitente que começou já na década de 1830 e não culminou até o final do século, com a (re)invenção do conceito de utilidade marginal por Jevons, Walras e Menger no início da década de 1870 marcando um marco importante nesse processo, Blaug (1973, pp. 6–7, 11) e Coats (1973b, p. 337). Como Meek (1973a, pp. 243–4) coloca:

\begin{quote}
O termo ‘revolução’ aqui é um tanto impróprio. A mudança na atmosfera geral foi real o suficiente, mas as ideias principais dos ‘revolucionários’ não eram de modo algum tão novas quanto às vezes eles gostam de afirmar. Muitas dessas ideias já haviam sido apresentadas – muitas vezes em uma forma surpreendentemente ‘avançada’ – nos anos anteriores a 1870, particularmente no curso dos debates sobre a teoria ricardiana que ocorreram nas décadas de 1820 e 1830.
\end{quote}

Apesar de suas muitas diferenças, os marginalistas compartilhavam muitas coisas em comum. Aqui nos concentramos principalmente nos escritos de Stanley Jevons e Léon Walras que, apesar de suas diferenças, compartilhavam mais entre si do que com Menger. Como diz Walras (1954, p. 36), ‘O trabalho do Sr. Jevons e o meu, longe de serem mutuamente competitivos em qualquer sentido prejudicial, realmente se apoiam, completam e reforçam um ao outro em um grau singular’. Carl Menger será tratado com mais detalhes na Seção 3 abaixo, no contexto do \textit{Methodenstreit}, e no Capítulo 13.

Um ponto de partida comum dos marginalistas foi a rejeição do que Jevons (1957 [1871], p. li) chamou de ‘as principais doutrinas da economia Ricardo-Mill’, como a teoria do fundo salarial, a teoria do valor baseada no custo de produção e a taxa natural de salários, pp. xlv–xlvi. Ele se refere rudemente a Ricardo como, ‘aquele homem capaz, mas equivocado, [que] desviou o carro da ciência econômica para uma linha errada’, p. xvi. ‘A única esperança de alcançar um verdadeiro sistema de Economia’, diz Jevons, p. xlv, ‘é descartar, de uma vez por todas, as suposições labirínticas e absurdas da escola ricardiana’. A única maneira de fazer isso, de acordo com Menger, é através de uma reforma completa da economia política. Tal reforma era necessária por causa do que ele considera terem sido as inadequações da Escola Clássica, que a tornaram incapaz de resolver ‘o problema de uma ciência das leis da economia nacional satisfatoriamente’, Menger (1985 [1883], p. 29). ‘Mesmo antes do aparecimento da escola histórica de economistas alemães’, escreve ele, ‘cresceu cada vez mais a convicção de que a crença anteriormente prevalecente na perfeição de nossa ciência era falsa e que, ao contrário, a ciência precisava de uma revisão completa’, pp. 27–8. A reforma e a revisão da ciência econômica tornaram-se o \textit{leitmotif} dos marginalistas. Para Menger (1985, p. 28), existem três caminhos abertos para a reforma da economia política. ‘Ou uma reforma da economia política teria que ser tentada com base nas visões anteriores de sua natureza e problemas e na doutrina fundada por Adam Smith... ou então novos caminhos teriam que ser abertos para a pesquisa’, p. 28. Uma dessas novas tentativas de reforma foi entregue pelos historicistas alemães, mas não para a satisfação de Menger, veja o Capítulo 5 e a Seção 3 abaixo. Os marginalistas ofereceram uma opção diferente.

Apesar da aversão dos marginalistas às doutrinas clássicas, dois dos pilares sobre os quais se assentou a revolução marginalista foram, na verdade, transmitidos por escritores clássicos. Primeiro está o método dedutivo abstrato de Ricardo, Senior, Mill e Cairnes, e segundo é o utilitarismo de Jeremy Bentham. Um dos objetivos básicos dos marginalistas era a transformação da economia de uma arte em uma ciência pura. ‘A principal preocupação do economista’, diz Walras (1954, p. 52), ‘é perseguir e dominar verdades puramente científicas’. Para fazer isso, a ciência econômica deveria se livrar de ‘vestígios e sobrevivências pré-científicos’, Winch (1973, p. 60), tais como os elementos sociais, históricos, filosóficos, políticos e éticos. Ao contrário da economia política clássica (e da Escola Histórica), ela deve se tornar uma ciência livre de valores em busca da verdade pura e desprovida de questões normativas sobre o que deve ser feito. Em outras palavras, a ciência pura deve ser claramente distinguida do que Walras, seguindo Senior e Mill, chama de artes ou ciências aplicadas e da ética ou ciências morais. A economia aplicada, segundo Walras (1954, p. 60), trata da questão do ‘que deve ser feito do ponto de vista do bem-estar natural’, deixando a ética ou as ciências morais com o que ‘deve ser feito do ponto de vista da justiça’.

Para ilustrar este ponto, Walras retoma a definição de economia política de Smith. Para começar, ele elogia Smith por ser o autor da primeira tentativa bem-sucedida, ‘de organizar o objeto da economia política como um ramo distinto de estudo’, p. 51. O que ele acha insatisfatório, no entanto, é a maneira como Smith define seu objeto de estudo que, na visão de Walras, levou sua análise na direção errada. De acordo com Smith, a economia política consiste em dois objetos: ‘primeiro, fornecer subsistência para o povo... e, segundo, fornecer ao estado ou comunidade receita suficiente para os serviços públicos’, citado em Walras (1954, p. 52). De acordo com Walras, no entanto, isso é exatamente o que a economia deveria evitar para se tornar uma ciência pura. ‘A característica distintiva da ciência pura’, diz ele, ‘é a completa indiferença às consequências, boas ou más, com que prossegue na busca da verdade pura’. O que Smith está descrevendo com sua definição de economia política não é a ciência da economia pura, mas a arte da economia aplicada, que trata da produção de riqueza social e da melhoria do bem-estar individual, pp. 54, 60. As ciências aplicadas devem ser mantidas estritamente separadas da economia pura, assim como a economia social (uma ciência moral ou ética), que trata de questões de propriedade, justiça e distribuição, pp. 76–80. As artes, diz Walras, pp. 61–2, ‘aconselham, prescrevem e dirigem’, enquanto a ciência ‘observa, descreve e explica’. As ciências aplicadas e morais lidam com fenômenos humanos (‘a operação da vontade humana’), enquanto ‘as operações das forças da natureza constituem o objeto do que é chamado ciência natural pura ou ciência’, p. 61.

Assim, Walras identifica ciência pura com ciência natural. Se a economia deve se tornar uma ciência pura, ela deve ser construída como se fosse uma ciência natural. Ele deixou claras suas intenções ao adicionar o adjetivo ‘pura’ no próprio título de seu livro, \textit{Elementos de Economia Pura}. Para Walras, a economia pode se tornar uma ciência pura ao desviar a atenção dos processos de crescimento e distribuição para o processo de troca e a determinação de preços – outra transformação básica trazida pela revolução marginalista. ‘A economia pura é, em essência, a teoria da determinação de preços sob o regime hipotético de concorrência perfeitamente livre’, diz Walras, p. 40. Com o marginalismo, a teoria da troca não apenas se torna o foco das atenções, mas também o próprio modelo de toda a ciência econômica, p. 44:

\begin{quote}
A teoria da troca baseada na proporcionalidade de preços às intensidades das últimas necessidades satisfeitas (i.e. aos Graus Finais de Utilidade...) que foi desenvolvida quase simultaneamente por Jevons, Menger e eu mesmo... constitui o próprio fundamento de todo o edifício da economia.
\end{quote}

Como, porém, todas as coisas obtêm seu valor da escassez, o valor de troca ‘participa do caráter de um fenômeno natural’, p. 69. Uma coisa tem que ser escassa para ter qualquer valor de troca. Assim, o problema da escassez e da escolha está sendo introduzido como parte, se não a parte principal, do objeto da ciência econômica. Não apenas isso: como o valor de troca é uma magnitude que é mensurável, e o objeto da matemática é estudar magnitudes, segue-se naturalmente que ‘a teoria do valor de troca é um ramo da matemática’, p. 70. Para Jevons (1957, p. vii, xxi), como a economia ‘trata o tempo todo com quantidades, deve ser uma ciência matemática’, tanto que ‘todos os escritores econômicos devem ser matemáticos na medida em que são científicos de todo’, veja também a citação de abertura de Walras. Para que a economia se torne uma ciência pura, então, ela deve ser tratada como uma ciência natural e, como lida com quantidades, deve tornar-se um ramo da matemática. Admitindo-se isso, torna-se óbvio que toda a economia política clássica e a Escola Histórica caem automaticamente fora do reino da ciência pura e tornam-se ou uma arte, uma ciência moral ou ambas. Assim, com os marginalistas, a busca por uma ciência econômica livre de valores em pé de igualdade com as ciências naturais é iniciada; esta jornada foi, até certo ponto, iniciada por Senior, Mill e Cairnes, e foi consolidada com os tratados metodológicos de John Neville Keynes em 1890 e Lionel Robbins em 1932, veja Capítulos 7 e 12, respectivamente.

Junto com essa mudança na definição do que constitui uma ciência econômica pura, e a mudança de ênfase que a acompanhou de questões macro-dinâmicas para micro-estáticas, os marginalistas adotaram o método abstrato/dedutivo, que está associado à outra grande divisão no pensamento econômico (ver Capítulo 2), em seu uso putativo como o método científico exclusivo de investigação econômica. Nisso, mais uma vez, eles seguem os passos de Ricardo, Mill e Cairnes. ‘Eu acho’, diz Jevons (1957, pp. 16, 17), ‘que John Stuart Mill está substancialmente correto ao considerar nossa ciência como um caso do que ele chama de método Dedutivo Físico ou Concreto’; esta visão é ‘quase idêntica à adotada pelo falecido Professor Cairnes’. Muito parecido com Mill, no entanto, Jevons qualifica seu uso do método dedutivo invocando a necessidade de verificação através do método empírico. Apenas tal verificação é extremamente difícil devido à falta de dados apropriados, pp. 21, 22:

\begin{quote}
Não hesito em dizer... que a Economia poderia ser gradualmente erguida em uma ciência exata, se apenas as estatísticas comerciais fossem muito mais completas e precisas do que são atualmente... O método dedutivo da Economia deve ser verificado e tornado útil pela ciência puramente empírica da Estatística... Mas as dificuldades desta união são imensamente grandes.
\end{quote}

Essas dificuldades, no entanto, não diminuem o valor da teoria dedutiva \textit{per se}. Jevons dá o exemplo do livre comércio para ilustrar este ponto. Embora, diz ele, os benefícios do livre comércio para todas as partes envolvidas não possam ser provados \textit{a posteriori}, através do uso de dados empíricos, isso não os priva de sua validade. Eles são verdadeiros apenas por virtude de serem provados logicamente através do método dedutivo. Como Jevons coloca, ‘eles devem ser acreditados porque o raciocínio dedutivo a partir de premissas de verdades quase certas nos leva confiantemente a esperar tais resultados’, p. 19. Walras (1954, pp. 71–2), por outro lado, defende o método dedutivo puro e simples. Nenhum recurso à verificação empírica é necessário. Todas as verdadeiras ciências ‘físico-matemáticas’, incluindo a economia, diz ele, ‘abstraem conceitos de tipo ideal que definem, e então, com base nessas definições, constroem \textit{a priori} todo o arcabouço de seus teoremas e provas. Depois disso, voltam à experiência não para confirmar, mas para aplicar suas conclusões’, p. 71.

Jevons simplesmente via a economia como um ramo de aplicação de leis científicas gerais. Schabas (1995, p. 198) reforça o ponto para Jevons citando uma palestra de 1876 sobre ‘O Futuro da Economia Política’:

\begin{quote}
Os primeiros princípios da economia política são tão amplamente verdadeiros e aplicáveis, que podem ser considerados universalmente verdadeiros no que diz respeito à natureza humana... Eu não perderia a esperança de traçar a ação dos postulados da economia política entre algumas das classes mais inteligentes de animais.
\end{quote}

Assim, a economia não é apenas universal para Jevons como uma expressão da natureza humana, ela também deve ser baseada na natureza animalesca dos seres humanos. Isso levou Leslie (1879a, p. 3), principal representante da Escola Histórica Britânica de economia, a satirizar:

\begin{quote}
O Sr. Jevons, embora favoravelmente disposto pela cultura filosófica e gostos em relação à investigação histórica na economia, instou, em nome da dedução do princípio aquisitivo, que mesmo os animais inferiores agem por um motivo semelhante, ‘como você descobrirá se interferir entre um cão e seu osso’. Um osso representa suficientemente bem o tipo de riqueza cobiçada por um cão, que tem um sistema cerebral comparativamente simples e poucos outros objetos. No entanto, você não pode prever a conduta nem mesmo de um cão por seu amor aos ossos, ou nenhum restaria no açougue. O cão tem consideração por seu dono e medo da polícia, e tem outras atividades.
\end{quote}

Em outras palavras, Jevons é inadequado como um teórico do cão e seu osso, quanto mais dos humanos e da economia como se fossem semelhantes ao cão e ao osso.

Como sugerido anteriormente, a visão macro-dinâmica da economia defendida pela economia política clássica dá lugar à análise de equilíbrio estático. O veículo básico nesta transformação, além do uso exclusivo do método dedutivo, foi o conceito de utilidade marginal, que se tornou a pedra angular sobre a qual todo o edifício neoclássico foi erguido. Como Schumpeter (1967 [1912], p. 181) coloca, ‘o conceito de utilidade marginal foi o novo fermento que mudou a estrutura interna da teoria moderna em algo bastante diferente daquela dos economistas clássicos’. Uma das principais implicações da adoção deste conceito é que o objeto da ciência econômica mudou da investigação das causas da riqueza e sua distribuição para o exame do comportamento econômico dos indivíduos, especialmente na forma do princípio da maximização (da utilidade). ‘Tentei tratar a Economia como um Cálculo de Prazer e Dor’, diz Jevons (1957, p. vi), no estilo típico benthamita e utilitarista. Dessa forma, a teoria objetiva do valor da economia política clássica (que assumiu a forma da teoria do valor-trabalho ou custo de produção) dá lugar a uma teoria subjetiva baseada na maximização da utilidade individual. De fato, tal foi a importância atribuída ao conceito de utilidade marginal por Jevons que, segundo Jaffé (1976, p. 517), ele ‘olhou para seu coeficiente diferencial como a arma letal com a qual golpearia para sempre a teoria clássica’. E para Schumpeter (1967, p. 190), ‘a teoria da utilidade marginal aceita o valor de uso como um fato da psicologia individual’. É importante, no entanto, estar ciente de que o uso da ‘subjetividade’, assim como a ‘psicologia individual’ mais amplamente baseada de Schumpeter, como base do valor, está aberto a mudanças de interpretação e conteúdo, e não está de modo algum confinado à familiar função de utilidade dada dos neoclássicos de hoje. De fato, o processo de redução da subjetividade e da utilidade a preocupações tão estreitas é parte integrante da formação do \textit{mainstream} em sua forma atual.

Assim, a adoção da dedução pura como um guia metodológico básico para uma economia pura foi acompanhada por outra ruptura metodológica. A preocupação com agregados econômicos, como classes ou a economia nacional, dá lugar ao que Menger chamou de ‘atomismo’ – o que mais tarde veio a ser conhecido como individualismo metodológico. Como visto no Capítulo 2, este último refere-se ao método de explicação pelo qual o todo é explicado em termos das propriedades de suas partes individuais (membros). Para Menger (1985 [1883], p. 93), apenas os indivíduos têm interesses. Nenhum corpo coletivo, como a economia nacional, pode ser analisado por si só sem referência aos seus membros individuais constituintes. Nas mãos dos marginalistas, e mais tarde com os escritores neoclássicos, o individualismo metodológico assume a forma específica do que tem sido chamado de ‘individualismo psicológico’, no qual o indivíduo é simplesmente considerado como um agente racional movido por algum motivo psicológico (utilitarista) para maximizar o próprio benefício, Zouboulakis (2002, p. 30).

O conceito de utilidade marginal também se provou instrumental em outros dois aspectos. Primeiro, sendo um conceito matemático, deu um grande ímpeto à matematização da economia. Como Hutchison (1953, p. 16) apropriadamente colocou ao extremo, ‘o que era importante na utilidade marginal era o adjetivo, mais do que o substantivo’. Segundo, forneceu uma justificativa para estreitar o escopo da investigação econômica para o estudo do problema da alocação sob escassez e a determinação de preços, focando nas relações de mercado tratadas isoladamente de seu contexto social e histórico. ‘Assim, originou-se uma economia diferente e muito mais “pura”’, Schumpeter (1967, p. 188), mas também ‘uma ciência muito contraída’, Black (1973, p. 106). Daí a mudança de economia política para economia, simbolicamente refletida no título da clássica \textit{magnum opus} da nova economia neoclássica, \textit{Principles of Economics}, de autoria de Alfred Marshall, mas veja o Capítulo 7. Com a excisão do político, também se foram o social e o histórico. Paradoxalmente, este estreitamento de escopo permitiu a expansão de suas fronteiras historicamente a um grau quase ilimitado: ‘os primeiros princípios da Economia Política são tão amplamente verdadeiros e aplicáveis, que podem ser considerados universalmente verdadeiros no que diz respeito à natureza humana’, Jevons (1876, p. 624) e Zouboulakis (1997, pp. 16–17). Como Winch (1973, p. 69) coloca:

\begin{quote}
O mérito principal da ciência reconstruída foi que ela demonstrou tanto a unidade quanto a universalidade das leis da escolha em situações econômicas. Ao definir o problema econômico como o de alocação de meios escassos entre usos alternativos... os proponentes do marginalismo enfatizaram a aplicação universal das leis da escolha humana.
\end{quote}

O que Jevons e os outros marginalistas buscavam era uma teoria geral com aplicabilidade universal, válida em todos os sistemas sociais, de Vroey (1975, p. 426).

Até agora, concentramo-nos no que os marginalistas compartilhavam em comum e como eles afetaram conjuntamente o curso dos eventos na evolução das ideias econômicas. No entanto, existem também diferenças importantes que se tornaram cada vez mais o foco das atenções. Embora a literatura tenha se concentrado principalmente em Menger, que é o mais diverso dos três, veja a Seção 3 e o Capítulo 13, existem também diferenças importantes entre Jevons e Walras. Já encontramos uma metodológica. Embora Walras seja a favor do método abstrato/dedutivo puro e simples, sem qualquer recurso à verificação empírica, esta última é um suplemento importante ao raciocínio abstrato no ‘método indutivo completo’ de Jevons, Black (1973, p. 110). No entanto, também existem grandes diferenças no que diz respeito às suas análises substantivas.

Para Jevons, a importação do utilitarismo benthamita para sua teoria é o ponto de partida, bem como o elemento mais onipresente de toda a sua análise, p. 23:

\begin{quote}
A teoria que se segue é inteiramente baseada em um cálculo de prazer e dor; e o objetivo da Economia é maximizar a felicidade comprando o prazer, por assim dizer, ao menor custo de dor... Não hesito em aceitar a teoria Utilitarista da moral que sustenta o efeito sobre a felicidade da humanidade como o critério do que é certo e do que é errado.
\end{quote}

Seu objetivo principal era exatamente ‘usar o utilitarismo para construir uma ciência exata da teoria do valor de troca’, Jaffé (1976, p. 518). Como ele coloca, ‘o valor depende inteiramente da utilidade’, Jevons (1957, p. 1). Por isso, ele inicia sua análise com um exame da teoria do prazer e da dor e da teoria da utilidade, que ele considera não como um aspecto intrínseco de um objeto, mas como uma avaliação subjetiva do prazer obtido através do uso de um objeto por um agente econômico. Em seguida, ele passa a discutir a teoria da troca, examinando o caso de dois indivíduos trocando duas mercadorias. A conclusão básica a que chega, ‘a pedra angular de toda a Teoria da Troca e dos principais problemas da economia’, como diz, é a seguinte, p. 95:

\begin{quote}
A razão de troca de quaisquer duas mercadorias será o recíproco da razão dos graus finais de utilidade das quantidades da mercadoria disponíveis para consumo após a troca ser concluída.
\end{quote}

Então Jevons considera o trabalho usando o mesmo aparato utilitarista, enquanto sua teoria da renda segue de perto a teoria clássica (ricardiana) dos rendimentos decrescentes da margem extensiva de trazer terras cada vez mais inferiores para o cultivo.

Quanto a Walras, seu sistema de equilíbrio geral, que é sua principal contribuição para a ciência econômica, é uma espécie de protótipo de uma teoria abstrata universal. O objetivo de sua análise é semelhante ao de Jevons: encontrar as condições de equilíbrio na troca. Mas a rota que ele segue é a oposta, e os meios que emprega são diferentes. O que diferencia a análise de Walras da de Jevons é o seu foco de atenção, seu sistema de equilíbrio geral, que é retratado como um conjunto de mercados inter-relacionados e é expresso em forma matemática como um conjunto de equações simultâneas. Diante disso, seu principal objetivo é provar matematicamente que a obtenção do equilíbrio em todos os mercados é possível, resolvendo este conjunto de equações em relação aos preços e quantidades relevantes. No equilíbrio, duas condições principais devem ser satisfeitas: primeiro, todos os agentes econômicos maximizam sua utilidade; segundo, a quantidade agregada demandada em cada mercado é igual à quantidade agregada ofertada. Matematicamente, a condição de utilidade máxima é alcançada quando ‘as \textit{raretés} dessas mercadorias são proporcionais aos seus preços após a conclusão da troca’, Walras (1954, p. 45). Esta é a condição também alcançada por Jevons, mas através de uma rota inteiramente diferente. Como Jaffé (1976, p. 515) mostrou persuasivamente, a \textit{rareté} (o ‘grau final de utilidade’ ou utilidade marginal de Jevons), em vez de ser o ponto de partida para a elaboração de uma teoria do consumo, como na teoria de Jevons, é simplesmente usada por Walras para fechar seu sistema de equilíbrio geral previamente derivado. Apesar do uso de funções de utilidade para fechar seu sistema, Walras não era um utilitarista e, em contraste com Jevons, o nome de Bentham não aparece uma única vez em todo o seu livro, p. 518, nem mesmo em sua \textit{Correspondência}. Como Jaffé (1976, p. 518) coloca, ‘Walras peremptória e despreocupadamente... postulou uma teoria de utilidade marginal mensurável sem mais delong, com o único propósito de arredondar sua teoria cataláctica de determinação de preços previamente formulada’. E acrescenta, p. 522, ‘Léon Walras sentiu que sua construção de equações simultâneas globais unidas pelo princípio da utilidade marginal provou que a \textit{rareté} é a causa do valor’.

Assim, o sistema de Walras é ‘rigorosamente estático em caráter... [e] é aplicável apenas a um processo estacionário’, este último referindo-se a ‘um processo que na verdade não muda por iniciativa própria’, Schumpeter (1937, pp. 165–6). Com o sistema de equilíbrio geral de Walras, a ciência econômica atinge seu apogeu no que diz respeito à sua natureza estática e ao seu conteúdo abstrato e matemático, que é totalmente divorciado da realidade, satisfazendo assim seu ditame de que ‘o estudioso tem o direito de buscar a ciência por si mesma’, pp. 71–2. É um sistema teórico que é tanto ‘institucionalmente vazio’ quanto ‘institucionalmente neutro’, Kaufman (2007, p. 10):

\begin{quote}
É vazio porque as instituições ou não existem (e.g. o dinheiro não tem papel teórico, as firmas são conjuntos de produção determinados tecnologicamente) ou são passivas e fatores de fundo dados exogenamente... É também institucionalmente neutro no sentido de que os resultados previstos são independentes tanto das atribuições de direitos de propriedade... quanto da forma de propriedade...
\end{quote}

Instituições e, deve-se acrescentar, história, simplesmente não importam no mundo walrasiano sem fricção e perfeitamente competitivo.

Apesar de sua grande ênfase na teorização pura, tanto Jevons quanto Walras também demonstraram interesse em assuntos mais práticos e questões de política. Apenas, para eles, estas últimas, que tratavam de questões normativas, são uma tarefa separada a ser totalmente distinguida da busca da ciência pura que lida com questões positivas. Em questões de política, tanto Jevons quanto Walras foram apoiadores de uma combinação de \textit{laissez-faire} com reformas moderadas na forma de legislação para questões como trabalho infantil e condições de saúde e segurança nas fábricas (Jevons) ou tributação da terra e nacionalização de monopólios naturais (Walras). Walras, este ‘reformador radical francês racionalista e otimista’ como Hutchison (1953, p. 216) o descreveu, chegou a se chamar de socialista, mas em essência sua posição política era uma ‘mistura de liberalismo tradicional e a doutrina da intervenção estatal’, Screpanti e Zamagni (1993, p. 170) e Oser e Blanchfield (1975, pp. 233–4).

\section{Carl Menger e o \textit{Methodenstreit}}

Em contraste com as escaramuças na Grã-Bretanha e a resistência a Jevons, veja o Capítulo 7, a controvérsia em torno do marginalismo na Alemanha foi intensa. A presença e a força da Escola Histórica inevitavelmente a levaram ao conflito com o marginalismo, até porque este último adotou o dedutivismo ricardiano que já havia sido rejeitado pela Escola Histórica. O conflito resultante atingiu seu clímax no \textit{Methodenstreit}, ou a ‘Batalha dos Métodos’, que ocorreu em 1883–4 entre Carl Menger e Gustav Schmoller. Então, o que essa batalha envolveu? Tendo introduzido os princípios gerais do marginalismo através do trabalho de Jevons e Walras, e tendo examinado de perto a Escola Histórica e Schmoller no capítulo anterior, o próximo passo é introduzir o outro protagonista do debate, Carl Menger, localizando-o tanto dentro do marginalismo quanto em relação à Escola Histórica. Menger é uma figura interessante e fundamental em muitos aspectos. Ele é considerado por muitos como o ‘estranho no ninho’ da revolução marginalista, Blaug (1997, p. 290). Ele é geralmente lembrado como um dos originadores do conceito de utilidade marginal e como o iniciador do \textit{Methodenstreit}, Schneider (1985, p. 1). Ao mesmo tempo, porém, ele também é considerado o pai da Escola Austríaca que, em sua versão neo-austríaca (representada principalmente por Mises e Hayek), apesar de seu forte subjetivismo, é totalmente oposta tanto ao marginalismo quanto ao neoclassicismo. O desenvolvimento da Escola Austríaca será escrutinado no Capítulo 13. Aqui nos concentramos em Menger e seu papel no \textit{Methodenstreit}.

Os \textit{Princípios de Economia} de Menger são seu único livro sobre teoria econômica e são considerados o \textit{locus classicus} da Escola Austríaca de economia, Streissler (1990, p. 33n). Nesta obra, o objetivo básico de Menger era fornecer os fundamentos científicos para a teoria econômica. Para ele, o objeto básico da teoria econômica é a investigação das ‘atividades práticas de homens que economizam’, Menger (2004 [1871], p. 48). Daí que o foco da teorização econômica deva ser ‘as condições sob as quais os homens se engajam em atividades providentes direcionadas à satisfação de suas necessidades’. Além disso, isso envolve:

\begin{quote}
Se e sob quais condições uma coisa é \textit{útil} para mim... se e sob quais condições ela é um bem econômico, se e sob quais condições ela possui \textit{valor} para mim e quão grande é a \textit{medida} do valor para mim, se e sob quais condições uma \textit{troca} econômica ocorrerá entre dois agentes que economizam e os limites dentro dos quais um \textit{preço} pode ser estabelecido se uma troca ocorrer.
\end{quote}

Menger queria descobrir as leis que governam os fenômenos econômicos que, segundo ele, e muito parecido com os outros dois marginalistas, assemelham-se às leis da natureza e são independentes da vontade humana, p. 48. Os princípios básicos sobre os quais seu trabalho repousa são, primeiro, o que Menger (enganosamente, pela terminologia de hoje) chama de ‘método empírico’, mas o que ele mais tarde chamaria de forma mais convencional de ‘método analítico-compositivo’. ‘No que se segue, tentei reduzir os fenômenos complexos da atividade econômica humana aos elementos mais simples que ainda podem ser submetidos à observação precisa... e... investigar a maneira pela qual os fenômenos mais complexos evoluem a partir desses elementos de acordo com princípios definidos’, Menger (2004, pp. 47–8). Segundo, esses ‘elementos mais simples’ nada mais são do que o ‘indivíduo que economiza’, de acordo com o princípio do individualismo metodológico. Como Ikeda (2006, p. 7) argumenta, Menger estava comprometido com o individualismo metodológico desde um estágio inicial. De fato, esta é a base para uma crítica contundente tanto a Adam Smith quanto à Escola Histórica, não menos por acreditarem que existe uma economia nacional independente das ‘economias individuais’ que a constituem:

\begin{quote}
Adam Smith e sua escola negligenciaram reduzir os complicados fenômenos da economia humana em geral, e em particular de sua forma social, a economia nacional, aos esforços das economias individuais, como seria de acordo com o real estado das coisas. Eles negligenciaram nos ensinar a compreendê-los teoricamente como o resultado de esforços individuais. Seus esforços visaram, antes, e, com certeza, subconscientemente na maior parte, nos fazer compreendê-los teoricamente do ponto de vista da ficção da economia nacional. Por outro lado, a escola histórica de economistas alemães segue essa concepção errônea conscientemente. É claro, no entanto, que sob a influência da ficção aqui discutida, um entendimento teórico dos fenômenos da economia nacional adequado à realidade não é alcançável. Além disso, o pouco valor das teorias econômicas prevalecentes encontra sua explicação em grande medida na visão básica errônea acima sobre a natureza da forma social atual da economia humana.
\end{quote}

Isso não é negar a existência da economia nacional, apenas que, em uma notável inversão da economia política marxista, a economia nacional é vista como uma ilusão (real) sustentada por economias individuais, apesar das aparências em contrário.

Por último, e como corolário, está o princípio do subjetivismo de Menger, pelo qual se entende que, ‘as valorações subjetivas colocadas pelos indivíduos em coisas que eles acreditam que satisfarão suas necessidades são a origem de toda atividade econômica’, Caldwell (2004, p. 25). Essas valorações subjetivas também se referem ao significado que os atores individuais anexam às suas ações, uma noção que foi desenvolvida mais tarde por Mises (1996) em seu conceito de ação humana, Capítulo 13, e por Max Weber através de sua concepção de compreensão (\textit{Verstehen}), Capítulo 11. É assim que Hayek (1973, p. 8) coloca a questão:

\begin{quote}
Menger acredita que ao observar as ações de outras pessoas somos assistidos por uma capacidade de \textit{compreender} o significado de tais ações de uma maneira que não podemos compreender eventos físicos. Isso está intimamente ligado a um dos sentidos em que pelo menos os seguidores de Menger falavam do caráter ‘subjetivo’ de suas teorias, pelo qual queriam dizer, entre outras coisas, que elas eram baseadas em nossa capacidade de apreender o significado pretendido das ações observadas. ‘Observação’, como Menger usa o termo, tem assim um significado que os behavioristas modernos não aceitariam; e implica um \textit{Verstehen} (‘compreensão’) no sentido em que Max Weber mais tarde desenvolveu o conceito.
\end{quote}

Para Menger, então, os indivíduos tanto fazem quanto \textit{interpretam} ações.

Com base nestes princípios, os temas básicos analisados nos \textit{Princípios} são, primeiro, a teoria subjetivista do valor em contraste com as teorias do valor baseadas no custo de produção ou no trabalho da escola clássica. Segundo, o \textit{processo} de formação de preços em oposição à determinação de preços walrasiana, Moss (1978, p. 17). Se, para Jevons e Walras, o problema era encontrar as propriedades dos preços de equilíbrio na troca, Menger estava mais interessado em explicar ‘o processo pelo qual esse conjunto de preços é atingido e muda’, em outras palavras, o processo de formação de preços, Wagner (1978, p. 66). Terceiro, o surgimento de instituições, que são consideradas as consequências não pretendidas da ação individual que economiza, e seu papel no desenvolvimento econômico. No entanto, enquanto Menger está comprometido com o individualismo metodológico e com as instituições como resultado não planejado de ações individuais, ele se opõe firmemente a um paralelo entre a evolução da sociedade humana e dos organismos naturais, e entre a ciência social e a natural. Isso ocorre porque o desenvolvimento social reflete mecanismos naturais apenas parcialmente, porque a ação humana é intencional na busca de fins criativos, mesmo que estes possam não resultar conforme o pretendido, Ikeda (2006, pp. 8–9). A análise de Menger procede passo a passo, começando com a análise de bens, depois passando para bens econômicos (escassos), então em seu Capítulo 3 para a análise do valor, que oferece a base para sua teoria de troca e preço, antes de, no capítulo final, chegar à explicação do surgimento e natureza de uma instituição específica, o dinheiro.

As primeiras escaramuças entre Menger e a Escola Histórica Alemã apareceram imediatamente após a publicação de seus \textit{Princípios}. Isso aconteceu apesar do fato de Menger dedicar o livro a Roscher e elogiar a Escola Histórica por fornecer o terreno sobre o qual seu próprio trabalho repousaria, ‘a reforma dos princípios mais importantes de nossa ciência aqui tentada é, portanto, construída sobre uma base lançada por trabalhos anteriores que foram produzidos quase inteiramente pela indústria de estudiosos alemães’, Menger (2004, p. 49). Gustav Schmoller, no entanto, pensava o contrário. Em uma breve resenha dos \textit{Princípios} publicada em 1873, ele atacou o livro de Menger por ser unilateral, e seu método por ser ‘remanescente de Ricardo em vez das tendências reinantes hoje nos círculos científicos alemães’, Schmoller (2004 [1873], p. 407). O que Schmoller rejeita é o método abstrato de Menger, cujos resultados não são de grande utilidade, pois ‘equivalem a nada mais do que novas formulações de tópicos convencionais abstratos em vez de soluções reais de problemas reais’, p. 407. Ele também critica a tentativa malsucedida do autor de emular o método das ciências naturais e de descobrir leis sociais com afinidades com as leis naturais e independentes da vontade humana, p. 408:

\begin{quote}
A base psicológica da vida econômica não é mutável, de acordo com o povo e a época? O autor não está reativando a velha e tendenciosa ficção inglesa, a saber, que a vida econômica poderia ser devidamente derivada da força motriz básica constante do homem médio abstrato? Não estão todos os problemas econômicos tornando-se para ele meras e puramente considerações privadas? As ciências naturais fizeram sua pesquisa precisa com balanças e microscópios; as abordagens que correspondem a elas na economia são a histórica, a estatística, etc.; se as ciências naturais quisessem proceder como o Dr. Menger faz na economia, elas teriam que explorar abstratamente o conceito da célula, do elemento químico e similares e derivar seus argumentos a partir daí. Isso também tem seu valor e sua justificativa, mas não é tanto um método exato quanto uma especulação sobre conceitos.
\end{quote}

O prelúdio para o \textit{Methodenstreit} já havia sido sinalizado. O que o desencadeou, no entanto, foi a publicação dez anos depois, em 1883, do tratado de Menger (1985 [1883]), \textit{Investigações sobre o Método das Ciências Sociais com Referência Especial à Economia}. Nele, ele faz um apelo apaixonado pela reforma da economia política por causa do que ele considera serem as inadequações da escola clássica. Uma dessas novas tentativas de reforma foi entregue pelos historicistas alemães, mas não para a satisfação de Menger. De fato, ele vê a tarefa principal deste livro como sendo uma crítica das doutrinas (objetivos e métodos) da Escola Histórica, que era então a ortodoxia prevalecente nas universidades alemãs, p. 32:

\begin{quote}
Fui guiado pelo pensamento de tornar a pesquisa no campo da economia política na Alemanha ciente de suas tarefas reais novamente. Pensei em libertá-la dos aspectos unilaterais prejudiciais ao desenvolvimento de nossa ciência, de libertá-la de seu isolamento no movimento literário geral e, assim, preparar o caminho para a reforma da economia política em solo alemão, uma reforma que esta ciência tão urgentemente necessita à luz de seu estado insatisfatório.
\end{quote}

Paralelo a essa tarefa está a ‘exposição positiva da natureza da análise teórica’ de Menger, Hayek (1934, p. 79). De fato, diz Hayek, as \textit{Investigações} ‘fizeram mais do que qualquer outro livro individual para tornar claro o caráter peculiar do método científico nas ciências sociais’, p. 79. A última contribuição deste livro é a oportunidade que deu a Menger para ‘uma elucidação da origem e do caráter das instituições sociais’, p. 79.

É assim que Menger resume sua posição no primeiro livro de suas \textit{Investigações}: ‘no livro anterior estabelecemos a diferença essencial entre as ciências históricas, teóricas e práticas da economia. Apontamos particularmente os erros daqueles que veem uma ciência “histórica” da economia política’, p. 97. Talvez o elemento mais onipresente desta parte do livro seja o esforço de Menger para esclarecer a diferença entre teoria e história. Ele faz esforços contínuos e vigorosos para tornar o mais claras possíveis as distinções entre o individual e o geral, o histórico e o teórico, a ‘plena realidade empírica’ e o ‘conhecimento teórico exato’, etc. De acordo com Menger, existem duas maneiras diferentes de analisar os fenômenos sociais: a individual e a geral. A primeira refere-se a ‘fenômenos concretos em sua posição no tempo e no espaço e em suas relações concretas uns com os outros’, enquanto a segunda refere-se às ‘formas empíricas recorrentes na variação destes, cujo conhecimento forma o objeto de nosso interesse científico’, p. 35. Correspondendo a esses dois pontos de vista diferentes, existem dois tipos diferentes de conhecimento científico. Por um lado, está o que ele chama de ‘orientação realista-empírica da pesquisa teórica’, que se refere à investigação de fenômenos sociais em sua ‘“plena validade empírica”, isto é, na totalidade e em toda a complexidade de sua natureza’, p. 56. Por outro lado, existe o ‘conhecimento teórico exato’, referindo-se à ‘determinação de leis de fenômenos que... devem ser designadas pela expressão “leis exatas”’, p. 59. De fato, a busca por ‘leis exatas’ como a tarefa principal da ciência econômica tornou-se o \textit{leitmotif} do trabalho de Menger sobre metodologia. A análise histórica está mais próxima do primeiro tipo de conhecimento, pois se refere à investigação dos processos individuais de desenvolvimento de um fenômeno econômico, p. 43, enquanto a economia teórica está mais próxima do segundo tipo de conhecimento científico, pois trata da ‘natureza geral e conexão geral (leis) de fenômenos econômicos, p. 39. A história, diz Menger, examina ‘todos os lados de certos fenômenos’, enquanto as teorias exatas examinam ‘certos lados de todos os fenômenos’, p. 79.

Admitindo tudo isso, o principal problema com a economia alemã é ‘a confusão do entendimento histórico e teórico dos fenômenos sociais por um lado, e a concepção unilateral do problema teórico das ciências sociais como exclusivamente realista por outro’, pp. 74–5. Menger também reagiu fortemente aos aspectos coletivistas e organicistas da Escola Histórica, reafirmando firmemente o indivíduo como a unidade básica de análise, que é a base do que ele chamou de ‘atomismo’ (ou o que hoje em dia seria chamado de ‘individualismo metodológico’). Para Menger, apenas os indivíduos têm interesses. Nenhum corpo coletivo, como a economia nacional, pode ser analisado por si só sem referência aos seus membros individuais constituintes, p. 93. Pelo contrário, as coletividades, como instituições ou nações, são elas mesmas o resultado subjacente da ação individual, p. 93:

\begin{quote}
os fenômenos da ‘economia nacional’ não são de modo algum expressões diretas da vida de uma nação como tal ou resultados diretos de uma ‘nação econômica’... Pelo contrário, os fenômenos da ‘economia nacional’, tal como se apresentam a nós na realidade como resultados de esforços econômicos individuais, devem também ser interpretados teoricamente sob esta luz.
\end{quote}

A diferença com os métodos e motivos da Escola Histórica é clara.

Sobre a questão da motivação individual, Menger reconhece a existência de outros motivos, como ‘espírito público, amor ao próximo, costume, sentimento de justiça’, p. 84, bem como o autointeresse. O que o impele, no entanto, a basear suas investigações teóricas exclusivamente no motivo do autointeresse, é a necessidade de construir uma ciência exata da economia baseada na abstração e no método de isolamento. Como ele coloca, a teoria exata da economia política ‘não tem a tarefa de nos ensinar a compreender de forma geral e em sua totalidade os fenômenos sociais... Ela tem apenas a tarefa de nos proporcionar a compreensão de um lado especial da vida humana, com certeza o mais importante, o econômico’, p. 87. Tendo isolado a esfera econômica de outras esferas sociais, a tarefa da ‘orientação exata da pesquisa teórica’ é reduzir os ‘fenômenos humanos às expressões das forças e impulsos mais originais e gerais da natureza humana’, p. 86. De todos esses impulsos humanos, Menger escolhe a satisfação do bem-estar, que é a principal manifestação do ‘autointeresse humano’ na esfera econômica e material, e de longe a mais importante.

Com base na análise precedente, Menger identifica três grupos diferentes de ciências da economia: primeiro, ‘as ciências históricas e as estatísticas da economia’; segundo, ‘economia teórica’, ambas as quais já encontramos; terceiro, ‘as ciências práticas... com a tarefa de investigar e descrever os princípios básicos da ação adequada... no campo da economia nacional’, p. 38. Em sua definição de economia política, Menger inclui apenas as duas últimas, deixando a investigação histórica totalmente fora do escopo da economia política: ‘Por economia política’, diz ele, ‘entenderemos aquela totalidade das ciências teórico-práticas da economia nacional (economia teórica, política econômica e a ciência das finanças)’, pp. 39–40. Assim, deliberadamente, e possivelmente de forma provocativa, ele move sua posição o mais longe possível dos historicistas e exclui muito de seu esforço de estar dentro da órbita da economia política.

Em suma, Menger em 1883 reagiu às reivindicações indutivistas dos historicistas reafirmando firmemente o método dedutivo abstrato na investigação econômica e o indivíduo como a unidade básica de análise. Segundo ele, a abordagem dedutiva é o único método verdadeiramente científico. A indução pura baseada na mera observação não pode ser a base para a investigação científica. O indivíduo econômico baseado no princípio do autointeresse torna-se a pedra angular da análise econômica de Menger. O coletivismo metodológico de grande parte da economia política clássica e da Escola Histórica dá lugar ao individualismo metodológico, e os aspectos de economia do comportamento humano tornam-se o foco primário da investigação econômica de Menger. No entanto, uma das principais preocupações da economia política de Menger, perdida no despertar da revolução marginalista e suas consequências, foi a questão das origens e natureza das instituições sociais. Isso tornou-se vinculado a escolas de pensamento não ou mesmo anti-marginalistas (a Escola Histórica, o antigo institucionalismo, a Escola Neo-Austríaca). Por último, embora substantivamente Menger tenha pedido uma reforma da economia política longe das doutrinas da economia política clássica, metodologicamente e no que diz respeito à forma adequada de teorizar em assuntos econômicos, especialmente no que diz respeito ao uso da dedução, suas visões, assim como as dos outros marginalistas, eram muito próximas em alguns aspectos de escritores clássicos como Senior, Mill e Cairnes, Hutchison (1973b, p. 27).

A reação de Schmoller foi rápida e afiada. Em sua resenha das \textit{Investigações} de Menger, publicada no mesmo ano (\textit{Zur Methodologie der Staats- und Sozialwissenschaften}), seu ponto de partida, ou seja, a necessidade urgente de reforma da doutrina clássica, é o mesmo de Menger, mas por razões muito diferentes. ‘A medula de sua força’, diz ele, traduzido em Small (1924, p. 221), ‘secou, porque tentou compor seus resultados excessivamente em esquemas abstratos que careciam de toda realidade’. Assim, em vez de simplesmente mudar a substância da velha doutrina enquanto mantinha o método abstrato, como Menger desejava, o que era necessário era um começo totalmente novo, ‘olhando para as coisas de um novo ângulo’. Tal nova época, no entanto, não poderia vir através da reciclagem e do refinamento adicional de teorias cada vez mais abstratas, mas sim através da coleta de material cada vez mais descritivo, histórico e estatístico. Isso equivale a uma poderosa reafirmação da necessidade da indução como o veículo para a reforma da economia política (clássica), porque as teorias gerais, na medida em que existem, têm que ser firmemente fundamentadas na realidade. ‘A ciência descritiva’, diz ele, ‘fornece as preliminares para a teoria geral... Toda descrição completa é, portanto, uma contribuição para o estabelecimento do caráter geral da respectiva ciência’, p. 220.

A resposta de Menger em \textit{Os Erros da Escola Histórica} (\textit{Die Irrtümer des Historismus}), publicada em 1884, não cede terreno ao seu oponente. Ele essencialmente recapitula os argumentos de suas \textit{Investigações}, criticando os estudiosos alemães por sua ênfase unilateral em estudos históricos e estatísticos e em direção à ‘investigação particularista’, por sua falha em ‘fazer uma distinção nítida entre as divisões teóricas e práticas da ciência econômica’ e por subvalorizar a necessidade de investigação econômica teórica, traduzido em Small (1924, pp. 221–2). Mais uma vez, ele faz uma distinção nítida entre o método empírico e o ‘método de isolamento... [cujo] objetivo é nos permitir compreender os fenômenos sociais não “em sua plena validade empírica”, mas apenas na medida em que seu lado econômico é concernido’. E, embora reconheça a utilidade da ‘história da economia pública’, ele a considera não como parte da economia política, mas como uma auxiliar dela, pp. 225, 227, veja também abaixo.

Para começar, tanto Menger quanto Schmoller, e seus seguidores, ‘bateram chifres’ e mantiveram suas posições iniciais, Balabkins (1988, p. 45). No entanto, cerca de uma década depois, os rivais já começaram a suavizar suas posições um em relação ao outro. Assim, Schmoller em seu \textit{Grundriss} chamou a atenção para a importância da abstração para qualquer observação significativa, enquanto escrevia em 1911, aceitando a necessidade de um economista usar tanto a indução quanto a dedução, pp. 45, 57. Da mesma forma, para Menger escrevendo em 1894, citado em Hutchison (1973b, p. 35):

\begin{quote}
Se os economistas históricos alemães foram frequentemente descritos... como representantes do método indutivo, e os economistas austríacos do método dedutivo, isso não corresponde aos fatos... Ambos reconhecem que a base necessária para o estudo de fenômenos reais e suas leis é a da experiência. Ambos reconhecem, – posso muito bem assumir, – que a indução e a dedução estão intimamente relacionadas, apoiando-se mutuamente e sendo meios complementares de conhecimento.
\end{quote}

Böhm-Bawerk (1890, p. 249), por outro lado, um dos seguidores e discípulos imediatos de Menger, escrevendo em 1890, sugere que:

\begin{quote}
a questão não é se o método histórico ou o método ‘exato’ é o correto, mas apenas se, ao lado do inquestionavelmente justificado método histórico da economia, não reconheceremos o método ‘isolador’. Muitos, eu entre eles, sustentam que ele deve ser reconhecido.
\end{quote}

Um clima de reconciliação, então, logo tomou o lugar da postura resiliente anterior de ambos os lados, vindicando, do ponto de vista dos protagonistas deste debate, o ditame de Schumpeter a respeito das ‘energias desperdiçadas’ que esta batalha teórica acarretou.

\section{As consequências}

Como já foi argumentado, a crescente confiança no método dedutivo pelos marginalistas, em conjunto com a mudança no objeto da investigação econômica e uma ênfase correspondente na maximização individual e no princípio da utilidade, foram instrumentais no estreitamento do escopo da investigação econômica. Confinou a economia a problemas que poderiam ser resolvidos aplicando o processo de raciocínio lógico-matemático, afastando a ciência econômica do confronto imediato com o mundo real. Ao mesmo tempo, o uso da análise marginal (uma ferramenta essencialmente matemática) e o conceito de equilíbrio (emprestado da mecânica estática) tornaram a economia mais suscetível à análise matemática, empurrando a ciência econômica ainda mais pelo caminho da abstração e do formalismo, Deane (1978, pp. 83–6, 95), Pheby (1988, p. 18) e Mirowski (1989b).

A transformação no escopo e no método da investigação econômica significada pelo triunfo do marginalismo encapsulou um triplo reducionismo. Primeiro é um reducionismo individualista, através do qual agentes econômicos coletivos são substituídos por indivíduos como a unidade básica de análise, e a economia é tratada como a mera agregação de seus membros individuais. Segundo é um reducionismo associal, onde a economia é tratada isoladamente de seu contexto social mais amplo através da exclusão total de todas as relações sociais (além do mercado) da análise. Por último, há um reducionismo anti-historicista, através do qual a ciência econômica é totalmente divorciada da história, Screpanti e Zamagni (1993, p. 149). Uma consequência importante desta enorme transformação na ciência econômica é o surgimento das disciplinas de sociologia e história econômica, baseadas no divórcio prévio da economia da sociedade e da história, respectivamente.

A mudança da economia política para a economia anunciou a separação desta última das outras ciências sociais. O reducionismo antissocial e anti-historicista do marginalismo deu à ciência econômica uma justificativa para se desenvolver independentemente de outras disciplinas sociais, Deane (1978, p. 75). Como foi mostrado, isso não é simplesmente o resultado do triunfo do marginalismo. Em vez disso, foi o resultado de um processo longo e demorado, pelo qual cada um desses elementos (o coletivista/holista, o histórico e o social) foi gradualmente retirado da análise econômica. Assim, o elemento histórico já havia sido, em grande parte, excisado da análise dedutiva de Ricardo, Senior e Cairnes. O individualismo metodológico, como vimos, já estava parcialmente presente tanto em Smith quanto em Mill, e totalmente no utilitarismo de Bentham, enquanto – um ponto separado, mas crucial – o elemento social na forma de análise de classe foi o último a ser derrubado.

De fato, o maior obstáculo para este reducionismo em evolução tem sido e continua sendo a localização analítica da classe. Pois o capitalismo e seus teóricos fazem referência inequívoca e inevitável à classe – um atributo coletivo – através do reconhecimento e análise inevitáveis das categorias mais imediatas de salários, lucros e rendas. O deslocamento da análise de classe substantiva, e o movimento correspondente da economia política para a economia, precisaram esperar pelo triunfo do marginalismo no final do século XIX, que estabeleceu a grande maioria dos princípios da economia \textit{mainstream} como a conhecemos hoje. Significativamente, a ideia de classe ainda prevalece de forma diluída após a revolução marginalista, mesmo entre seus ardentes proponentes, não menos na teoria da distribuição. A ideia de que os salários, lucros ou rendas são simplesmente um preço como qualquer outro – que a distribuição é simplesmente um corolário da teoria dos preços e não sua pré-condição – provou ser difícil de aceitar, e este elemento da tradição clássica persistiu. Em certo sentido, o apego persistente à classe na economia marginalista, não menos através do uso de equilíbrio parcial por Marshall em oposição ao equilíbrio geral, representa um domínio precário sobre a especificidade histórica e social à medida que a economia política dá lugar à economia, Fine (1980c) e Capítulo 7.

Além do golpe mortal desferido na classe pelo individualismo metodológico, o que é distintivo no marginalismo e na economia neoclássica é, primeiro, a introdução dos conceitos de equilíbrio e utilidade marginal – até porque ambos são conceitos historicamente vazios – facilitando assim um conteúdo matemático. Segundo é a combinação de todas as três formas de reducionismo dentro do mesmo aparato analítico. Nesse sentido, o triunfo do marginalismo deu um impulso adicional e decisivo ao processo de separação da economia das outras ciências sociais, um processo que já havia começado antes de os marginalistas oferecerem seus tratados.

A economia deveria agora confinar-se à análise da economia, considerada simplesmente como o processo de mercado, focando nos aspectos econômicos do comportamento em abstração de quaisquer outras influências sociais. Espaço foi assim criado para o surgimento de outras ciências sociais, que preencheriam as lacunas induzidas pela dessocialização da economia. Duas características importantes são cruciais na remoção do social da economia. Primeiro, o foco da análise não apenas mudou para o indivíduo, mas também para a forma particular de comportamento otimizador através da maximização da utilidade (com a maximização do lucro pelas firmas ou empreendedores como corolário) ou individualismo psicológico de um tipo especial e limitado. Segundo, a economia tornou-se identificada com o mercado, com as relações sociais e políticas mais amplas desaparecendo no fundo dado exogenamente, a serem estudadas por outras ciências sociais. Em outras palavras, a economia como relações de mercado foi constituída como um objeto de estudo distinto, com a disciplina da economia para realizar a tarefa. Houve, como Hicks (1976) chamou, uma mudança de atenção da ‘plutologia’, ou a produção e crescimento da riqueza, para a ‘cataláctica’, ou a alocação de recursos, ela própria ultimamente reduzida a uma lógica de escolha individual, veja também Zafirovski (2001).

Não é por acaso que o estabelecimento da ortodoxia neoclássica, com sua visão estreita, associal e a-histórica da economia e do que constitui o objeto de sua investigação, foi logo seguido pelo surgimento da sociologia e da história econômica como disciplinas separadas. O nascimento da história econômica como uma disciplina separada será tratado no Capítulo 8, e Milonakis e Fine (no prelo), enquanto o surgimento da sociologia como um campo de estudo distinto será examinado no Capítulo 12.

Por enquanto, basta dizer que o nascimento da história econômica como uma disciplina separada pode ser atribuído a dois desenvolvimentos paralelos. Primeiro é o fracasso a longo prazo da Escola Histórica em estabelecer-se como uma escola alternativa de pensamento e teoria econômica; segundo, a separação da economia neoclássica da história através do estreitamento de seu escopo, desocupando espaço para o novo campo ou disciplina surgir. O episódio que serviu de catalisador neste processo, ao trazer à tona todas as principais diferenças entre o método histórico e a dedução abstrata, foi o \textit{Methodenstreit}.

Em suma, enquanto os economistas históricos tentavam tornar a história o objeto da investigação econômica, eles foram sitiados ou, mais exatamente, superados pela natureza totalmente a-histórica do marginalismo, baseada no uso do método dedutivo e no conceito de equilíbrio. Com o benefício da retrospectiva, com quaisquer que sejam os méritos intelectuais relativos, os marginalistas eventualmente conquistaram uma vitória decisiva no \textit{Methodenstreit}. Isso ocorre na prática, pelo menos no que diz respeito à disciplina da economia. A abordagem dos historicistas para a investigação econômica é geralmente considerada como tendo sofrido de duas fraquezas principais. Uma era sua aparente relutância em abstrair para isolar os fatores causais básicos na análise econômica. Como Schumpeter (1994, p. 812) descreve, falando de Schmoller em particular, ‘nada no cosmos ou caos social está realmente fora da economia schmolleriana’. Esta relutância é considerada relacionada a outra limitação, ainda mais séria, da economia histórica – sua aversão à teoria abstrata. Isso, segundo Schumpeter (1967, pp. 172–3), é ‘a causa científica mais importante’ da reação contra a Escola Histórica. Mesmo Schmoller, que aceitou explicitamente a necessidade de teorizar, falhou em fornecer um arcabouço teórico adequado. Em suma, o fracasso da economia histórica em estabelecer-se como uma escola alternativa na economia é atribuído principalmente à sua falta de teoria e coerência, Hodgson (2001, ch. 6).

Tal avaliação, no entanto, não representa todo o quadro e, como tal, não é inteiramente precisa. Primeiro, o conteúdo e a interpretação da(s) Escola(s) Histórica(s) foram fortemente influenciados por seu debate com a ‘economia’, com suas fraquezas exageradas e seus pontos fortes ignorados. Assim, por exemplo, como visto detalhadamente no Capítulo 5, a análise histórica não é e não pode ser inteiramente desprovida de teoria.

De acordo com Pearson (1999, p. 551):

\begin{quote}
membros praticantes da ‘jovem’ GHSE [Escola Histórica Alemã de Economia] não sentiam nenhuma obrigação particular de evitar a generalização. Alguns deles trabalharam em teoria relativamente ‘pura’, como Knap (teoria monetária e economia dos transportes), Bücher (organização industrial), Lujo Brentano (teoria do salário), e até o próprio Schmoller (tributação e seguros)... Na época em que a ‘mais jovem’ GHSE alcançou proeminência com Weber e Sombart, a associação da economia alemã com o obscurantismo curatorial era totalmente insustentável.
\end{quote}

Em espírito semelhante, Shionoya (2001a, p. 11) argumenta que:

\begin{quote}
mesmo para Schmoller, o método histórico não era simplesmente direcionado para a acumulação de historiografia e monografias históricas; em vez disso, visava reunir materiais para, em última análise, construir uma teoria mais ampla para o arcabouço institucional da economia e suas mudanças históricas.
\end{quote}

Ao longo do tempo, a(s) Escola(s) Histórica(s) não foram desprovidas de teoria nem anti-teoria. O que elas se opunham com muita força era o tipo universal de teoria válida em todos os lugares e para todas as épocas, do tipo praticado pelos marginalistas e mais tarde pela economia neoclássica, bem como o tipo especulativo de teoria abstrata que não está firmemente ancorada na realidade. Ao mesmo tempo, porém, a extensão em que usavam a teoria variava entre diferentes escritores da tradição, com uma tendência, pelo menos no que diz respeito ao ramo alemão, de tornar-se cada vez mais teórica. Esta tendência atingiu o seu clímax nas obras dos membros do ramo mais jovem (principalmente Weber e Sombart), que eram primariamente teóricos, veja o Capítulo 11.

Segundo, as verdadeiras questões por trás do \textit{Methodenstreit} não eram simplesmente metodológicas, mas também (e principalmente) epistemológicas e substantivas. Epistemologicamente, de acordo com Bostaph (1978), as diferenças entre Menger e Schmoller eram muito maiores do que eles próprios poderiam ter percebido. Tratá-las simplesmente como envolvendo a questão da indução versus dedução, como é geralmente feito, é perder o ponto do debate, pp. 14–5. Bostaph contrasta o nominalismo de Schmoller com a posição aristotélica de Menger. Para Schmoller como um nominalista, ‘a essência [de fenômenos específicos] deveria ser obtida por uma sumarização sobre entidades de todas as suas características, em vez da apreensão ou percepção de uma característica central e definidora’, p. 10. Isso envolve um tipo descritivo de causalidade onde, seguindo David Hume, ‘a explicação da relação causal [é percebida] como mera uniformidade na sucessão’, em vez de ‘quaisquer conexões intrínsecas ou necessárias que unem esses eventos’, p. 10. Como resultado, porque o contexto empírico é sempre diferente, a busca por conceitos universalmente aplicáveis e leis absolutas está fadada ao fracasso, p. 11. O essencialismo aristotélico de Menger, por outro lado, considera todos os fenômenos como tendo algumas propriedades essenciais que fundamentam sua natureza, e busca relações causais entre os elementos individuais desses fenômenos através do método de abstração e isolamento, pp. 11–12:

\begin{quote}
Menger buscou os ‘elementos mais simples’ de tudo o que é real, as essências, a natureza (\textit{das Wesen}) do real. Em sua abordagem exata, ele usou o processo de abstração a partir dos fenômenos individuais do mundo empírico para descobrir suas essências, isolá-las e então utilizar os ‘elementos mais simples’ assim obtidos para deduzir ‘como fenômenos mais complicados se desenvolvem a partir dos mais simples, em parte até mesmo de elementos não empíricos do mundo real’.
\end{quote}

Porque esses ‘elementos mais simples’, que, é claro, nada mais são do que as características essenciais dos indivíduos, são derivados da realidade empírica, o método de Menger, de acordo com Bostaph, não pode ser chamado de \textit{a priori}. Ao mesmo tempo, substantivamente, o \textit{Methodenstreit} referia-se a problemas reais: análise da utilidade e do preço através do método abstrato/dedutivo para os economistas neoclássicos, versus a questão da evolução institucional e do desenvolvimento das economias nacionais tratada através do método histórico e ético para os historicistas, Hutchison (1973, pp. 34–5) e Shionoya (2001a, p. 11). Como o próprio Menger coloca em seu tributo de obituário a Roscher em 1894, citado em Hutchison (1973b, p. 35):

\begin{quote}
As diferenças que surgiram entre a escola austríaca e alguns dos economistas históricos alemães não foram de modo algum de método no sentido estrito da palavra... O real fundamento das diferenças, que ainda não estão completamente superadas, entre as duas escolas é algo muito mais importante: relaciona-se à visão diferente em relação aos \textit{objetivos} da pesquisa, e sobre o conjunto de tarefas que uma ciência da economia tem que resolver.
\end{quote}

Portanto, algo mais, além do método, estava no coração do \textit{Methodenstreit}, como sugerido por um de seus próprios protagonistas.

Terceiro, e de grande importância no \textit{Methodenstreit} e em seu resultado, foi a postura política geral de cada um de seus participantes em relação à ordem estabelecida. Mises (2003, p. 15), o principal economista neo-austríaco e um defensor do liberalismo, em seu estilo característico, coloca este ponto em termos inequívocos. ‘A filosofia britânica do livre comércio triunfou no século XIX nos países da Europa Ocidental e Central’, diz ele. E continua:

\begin{quote}
Mas muito em breve o governo de Bismarck começou a inaugurar sua \textit{Sozialpolitik}, o sistema de medidas intervencionistas como legislação trabalhista, previdência social, atitudes pró-sindicatos, tributação progressiva, tarifas protetoras, cartéis e dumping. Se alguém tentar refutar a crítica devastadora feita pela economia contra a adequação de todos esses esquemas intervencionistas, será forçado a negar a própria existência – para não mencionar as reivindicações epistemológicas – de uma ciência da economia... É isso que todos os defensores do autoritarismo, da onipotência governamental e das políticas de ‘bem-estar’ sempre fizeram. Eles culpam a economia por ser ‘abstrata’ e defendem um modo ‘visualizador’ de lidar com os problemas envolvidos. Eles enfatizam que os assuntos neste campo são complicados demais para serem descritos em fórmulas e teoremas. Eles afirmam que as várias nações e raças são tão diferentes umas das outras que sua ação não pode ser compreendida por uma teoria uniforme... O significado político da Escola Histórica consistiu no fato de que ela tornou a Alemanha segura para as ideias cuja aceitação tornou populares entre o povo alemão todas aquelas políticas desastrosas que resultaram nas grandes catástrofes. Schmoller e seus amigos e discípulos defenderam o que foi chamado de socialismo de estado; i.e., um sistema de socialismo – planejamento – no qual a alta administração estaria nas mãos da aristocracia Junker. Era este tipo de socialismo que Bismarck e seus sucessores visavam.
\end{quote}

Em outras palavras, para Mises, a postura anti-teórica da Escola Histórica Alemã explica-se pelo fato de que a economia teórica da época fornecia o terreno científico para a legitimação das políticas de livre comércio e contra o intervencionismo que eles defendiam. O triunfo deste último na Alemanha de Bismarck oferece uma explicação para o sucesso que esta escola desfrutou na Alemanha durante a maior parte de um século. Sua estreita associação com o que Mises enganosamente chama de ‘socialismo de estado’, no entanto, também determinou seu destino no longo prazo. Schumpeter (1967, pp. 172–3) sugere isso quando escreve que esta última ‘tinha se associado a tendências políticas da mesma forma que os economistas clássicos fizeram em seu tempo. E, como estes últimos, agora tinham que pagar o preço por isso’.

Embora as opiniões de Schumpeter e Mises expressas acima devam ser vistas com cautela por causa de suas próprias fortes visões políticas liberais, que às vezes ‘nublaram sua percepção da história de seu assunto’, Streissler (1990, p. 40), elas certamente, neste caso, contêm um forte elemento de verdade. Os elementos políticos e ideológicos, como vimos, desempenharam um papel decisivo no desenvolvimento das ideias econômicas. O próprio Schmoller está absolutamente consciente deste fato. Como ele coloca, citado em Koslowski (2002, p. 151):

\begin{quote}
A economia política de hoje representa uma concepção ética e histórica do estado e da sociedade, em vez de uma determinada pelo realismo e materialismo. Da teoria pura do mercado e da troca, um tipo de ‘economia do nexo monetário’, que já foi uma arma de classe dos ricos, tornou-se mais uma vez uma grande ciência moral-política. Ela analisa não apenas a produção, mas também a distribuição de bens, os processos de agregação de valor, bem como as instituições econômicas, e coloca o homem em vez de bens e capital no centro do esforço científico.
\end{quote}

A atitude crítica de Schmoller em relação ao sistema capitalista (e a da Escola Histórica em geral), e a defesa da intervenção governamental e da reforma social, devem ser contrastadas com a abordagem geralmente panglossiana dos austríacos, onde, na verdadeira tradição liberal, a condição humana é apresentada como o melhor de todos os mundos possíveis, e onde a harmonia social e a eudaimonia são objetivos viáveis dentro do contexto dado. Como Mises (2003, p. 16) colocaria mais tarde, ‘cada nova geração adicionará algo ao bem realizado por seus antepassados. Assim, a humanidade está na véspera de um avanço contínuo em direção a condições mais satisfatórias. Progredir firmemente é a natureza do homem... O estado ideal da sociedade está diante de nós, não atrás de nós’. Dito isto, estudos recentes parecem estar divididos sobre a questão da predisposição política de Menger. Este quadro nebuloso é em parte resultado do fato de Menger ter escrito preciosamente pouco sobre questões de política. De acordo com uma visão, que parece predominar, as evidências existentes sugerem que ele foi muito provavelmente um ‘campeão intransigente do \textit{laissez-faire}’, Kirzner (1990, pp. 94–7) e Streissler (1990). Outros estudiosos, no entanto, consideram Menger mais como um social-democrata, que simpatizava com os fracos e os pobres à custa da aristocracia, e apontam para outras evidências que sugerem que, ao contrário da sabedoria convencional, Menger era um defensor da intervenção governamental, Wagner (1978, pp. 1–2) e Boehm, citado em Kirzner (1990, p. 94).

\section{Observações finais}

As causas do sucesso da revolução marginalista são multifacetadas. Proeminentes são os problemas enfrentados pela própria economia política clássica, que levaram à crise da economia ricardiana no final da década de 1860 e 1870. Segundo Hutchison (1978, p. 58), ‘no espaço de alguns anos, no final da década de 1860 e início de 1870, a estrutura clássica de “teoria” sofreu um colapso de confiança notavelmente súbito e rápido, considerando quão longa e autoritária foi sua dominância na Grã-Bretanha’. Esta crise refletiu-se no jantar realizado em 1876 para honrar o centésimo aniversário da publicação de \textit{A Riqueza das Nações}. Revelou que o consenso construído em torno da popularização do ricardianismo por Mill havia evaporado. Tornou-se aparente que havia divisões profundas entre os participantes em relação a cada aspecto da economia política (método, papel da história, intervenção estatal, distribuição, livre comércio, etc.). Isso levou Bagehot, citado em Hutchison (1953, p. 6), a comentar que a economia política ‘jaz morta na mente pública’. No mínimo, a crise serviu à função de limpar o terreno para o surgimento da nova ortodoxia – mas nada mais. A crise de um dado paradigma é uma coisa, e o que deve substituí-lo é outra bem diferente. E, como vimos, havia três principais contendores: o marxismo, o historicismo e o marginalismo, Hutchison (1953, pp. 1–6 e 1978, p. 1) e Koot (1987, pp. 10–14).

O sucesso do marginalismo também está associado à profissionalização da economia, que ocorreu no último quartel do século XIX e ajudou a estabelecê-la como uma ciência acadêmica. Esta profissionalização estava intimamente ligada à aceitação e consolidação adicional do marginalismo como base da ‘nova’ ciência econômica, que não ocorreu realmente até o final do século, Howey (1973, p. 25). ‘Por quase duas décadas’, diz Hutchison (1973a, p. 185), ‘houve na Grã-Bretanha um interregno um tanto confuso’. Outros marcos neste processo de consolidação foram a adoção do conceito por autores importantes como os seguidores austríacos de Menger, Wieser e Böhm-Bawerk, e o sueco Wicksteed na década de 1880. A economia neoclássica, no entanto, não atingiu proeminência até bem depois da publicação da \textit{magnum opus} de Marshall, \textit{Principles of Economics}, em 1890, Howey (1973, p. 25 e 1960, chs 26, 27), Blaug (1973, p. 14), Stigler (1973, pp. 310–20), Winch (1973, p. 60), Khalil (1987, p. 119), e Capítulo 7.

Sob esta luz, a Escola Histórica foi derrotada e, se sim, com quais efeitos imediatos e a longo prazo? Em retrospectiva, a resposta óbvia pareceria ser um afirmativo ressonante à derrota, em vista do triunfo inequívoco do marginalismo dedutivo, pelo menos dentro da economia. Mas o registro do próprio debate sugere o contrário, com os pontos reveladores e até valiosos da Escola Histórica sendo pelo menos parcialmente aceitos em princípio por seus oponentes (por exemplo, a necessidade de combinar dedução com indução), mas sendo cada vez mais ignorados na prática à medida que o marginalismo experimentava seus próprios imperativos e dinâmica. Nesse sentido, Jacob Viner (1991, pp. 238–9), citado em Hutchison (2000, p. 359), atinge o tom certo em sua eloquente avaliação do destino do \textit{Methodenstreit}:

\begin{quote}
Meus colegas teóricos me dizem que os teóricos conquistaram uma vitória definitiva nesta batalha quando Carl Menger, na década de 1880, demoliu Gustav Schmoller. Não posso concordar. Acredito que a batalha foi majoritariamente simulada, e que enquanto Schmoller certamente não levou louros, os que desde então foram concedidos a Menger por sua vitória nesta batalha são de papelão... O real desafio que Menger deveria ter enfrentado não era o de justificar em princípio o recurso à abstração pelos economistas, mas o de justificar a extensão e a maneira particular como ele e seus colegas teóricos a praticavam.
\end{quote}

Derrotada pelo marginalismo ou não, uma coisa é certa. Os efeitos da revolução marginalista, e do \textit{Methodenstreit} em particular, foram duradouros e de longo alcance. Um efeito particular para o qual o \textit{Methodenstreit} contribuiu, ainda que indiretamente, foi o nascimento da Escola Austríaca de economia, que se consolidou e tornou-se cada vez mais marginalista através da influência dos seguidores e discípulos de Menger, Böhm-Bawerk e Wieser. Mais tarde, porém, através da influência de Mises e Hayek, que deu à escola seu sabor neo-austríaco, o elemento marginalista foi perdido, embora ainda permanecesse cada vez mais individualista e subjetivista, veja o Capítulo 13.

Um segundo efeito foi o destacamento da economia teórica das questões de política. Ao contrário da Escola Histórica, cujo objetivo principal era construir uma economia histórica a serviço da reforma social para enfrentar os problemas prementes da época, enquanto enfatizava a dimensão ética da economia política, os marginalistas buscaram construir uma ciência econômica livre de valores baseada em leis exatas de maneira semelhante às ciências naturais. Esta tendência no pensamento econômico não era nova. Ricardo, Mill e Senior foram todos adeptos desta abordagem. O que o marginalismo fez foi ajudar a consolidar a distinção entre economia positiva e normativa, com a economia científica sendo confinada à primeira. Esta consolidação foi tornada explícita pelo tratado de John Neville Keynes sobre o método, \textit{The Scope and Method of Political Economy}, publicado em 1890, veja o próximo capítulo. O debate sobre a desejabilidade ou não de uma ciência econômica e social livre de valores recebeu um novo impulso através da publicação do ensaio de Weber sobre ‘A Objetividade nas Ciências Sociais e na Política Social’, onde ele se posicionou firmemente ao lado de uma ciência social livre de valores, veja o Capítulo 11, Seção 3. Este artigo deu origem a outro prolongado debate metodológico. Na economia, o apelo por uma ciência econômica positiva e livre de valores alcançou a posição suprema que continua a desfrutar até os dias atuais após a publicação do livro de Robbins, \textit{An Essay on the Nature and Significance of Economic Science}, publicado em 1932, mais de 40 anos depois, Balabkins (1988, pp. 63–4) e Capítulo 12. Assim, ocorreu uma mudança da economia política para a economia (positiva), proclamada pela primeira vez por Jevons e explicitamente retratada no título da \textit{magnum opus} de Marshall, \textit{Principles of Economics}.

Por último, e relacionado ao anterior, estava a separação da economia do social e do histórico. Construir uma ciência ‘exata’ da economia significava isolar a esfera econômica de outras esferas sociais. Este estreitamento do escopo da economia criou o espaço para o surgimento de outras ciências sociais, incluindo a sociologia, para lidar com os aspectos não racionais do comportamento humano, Capítulo 12, e a história econômica como, no máximo, uma auxiliar da economia, Capítulo 8. Assim, Max Weber, ele mesmo um membro da mais jovem Escola Histórica Alemã, que foi tão instrumental no nascimento da sociologia como disciplina, é também interessante como um autor que se encontrou preso entre os extremos da Escola Histórica e o recém-emergente marginalismo. Ele tentou reconciliar os dois lados através de uma convergência mútua. Sob esta luz, Weber pode ser visto como representando o ponto final da evolução e conciliação, no que diz respeito à Escola Histórica, e um representante proeminente da nova ordem ao buscar adicionar o social à racionalidade econômica. O resultado foi uma tentativa efêmera de economia social, para a qual voltamos no Capítulo 11, significando em última análise o desaparecimento tanto da Escola Histórica da qual derivou em parte, quanto também do ‘velho’ marginalismo de Marshall que esperava promover como um componente central, veja o próximo capítulo.

Como é aparente e geralmente aceito, o papel de Marshall na promoção do marginalismo é central. Isso não ocorre porque ele era obstinadamente a favor de seus princípios – muito pelo contrário. Pois, embora estivesse plenamente consciente do \textit{Methodenstreit}, ele prevaleceu no contexto de um meio intelectual diferente, embora um que incorporasse afinidades com a Escola Histórica Alemã em sua oposição ao marginalismo. Os principais adversários de Marshall encontravam-se entre a Escola Histórica Britânica, em relação à qual Marshall vacilou entre o compromisso e o conflito. De qualquer forma, o resultado foi testemunhar a marcha progressiva do marginalismo. Quaisquer que fossem suas próprias reservas, o objetivo de Marshall era estabelecer o marginalismo em princípio como o padrão central em uma economia profissionalizada. Ele próprio forneceu esse núcleo, temperado por concessões à importância de outros fatores, dando origem ao antigo marginalismo. Este último, porém, deu lugar ao novo marginalismo – a economia neoclássica como a conhecemos hoje – com a incômoda consciência intelectual fornecida pelas Escolas Históricas tendo desvanecido no esquecimento. Isso dentro da própria disciplina da economia. Enquanto o resíduo da Escola Histórica Alemã sobreviveu, ainda que intermitentemente, como economia social e institucionalismo americano, Capítulos 9–11, na Grã-Bretanha, a ideia de que a economia precisava de algo além do marginalismo encontrou expressão em outras formas. Deu origem à história econômica como uma disciplina separada (e como progênie da Escola Histórica Britânica), Capítulo 8, e à emergência da macroeconomia como distinta da microeconomia, Capítulo 14.

\end{document}